\documentclass{article}
\usepackage{amsmath,amssymb,amsthm}
\usepackage{algorithm}
\usepackage[noend]{algpseudocode}
\usepackage{enumitem}
\usepackage{hyperref}
\usepackage{geometry}
\geometry{margin=1in}
\newtheorem{theorem}{Theorem}
\newtheorem{lemma}{Lemma}
\newtheorem{assumption}{Assumption}
\newtheorem{definition}{Definition}
\newcommand{\E}{\mathbb{E}}
\newcommand{\R}{\mathbb{R}}
\newcommand{\norm}[1]{\left\|#1\right\|}
\newcommand{\ip}[2]{\left\langle #1,#2\right\rangle}
\title{AMSGrad: Algorithmic Description and Convergence Proof for Non‑Convex Smooth Objectives}
\author{}
\date{}
\begin{document}
\maketitle

\section*{Algorithm Name}
\textbf{AMSGrad} – Adam with monotone second‑moment estimates.

\section*{Mathematical Formulation}
Let $f:\R^{d}\rightarrow\R$ be the objective function.  
For iteration $t=1,2,\dots,T$ the algorithm maintains:

\begin{align}
g_t &:= \nabla f(w_t) \quad\text{(true gradient, or unbiased stochastic estimate)}\label{eq:grad}\\
m_t &= \beta_1 m_{t-1} + (1-\beta_1) g_t , \qquad m_0 =0,\label{eq:m}\\
v_t &= \beta_2 v_{t-1} + (1-\beta_2) g_t^{\odot 2}, \qquad v_0 =0,\label{eq:v}\\
\hat{v}_t &= \max (\hat{v}_{t-1},\, v_t),\qquad \hat{v}_0 =0,\label{eq:vmax}\\
\hat{m}_t &= \frac{m_t}{1-\beta_1^{t}},\qquad 
\hat{v}_t^{\text{bias}} = \frac{\hat{v}_t}{1-\beta_2^{t}},\label{eq:bias}\\
w_{t+1} &= w_t - \eta \,\frac{\hat{m}_t}{\sqrt{\hat{v}_t^{\text{bias}}}+\epsilon},\label{eq:update}
\end{align}
where $\odot$ denotes element‑wise product, $\beta_1,\beta_2\in[0,1)$ are exponential decay rates, $\eta>0$ the stepsize, and $\epsilon>0$ a small constant for numerical stability.

\section*{Pseudocode}
\begin{algorithm}[H]
\caption{AMSGrad}
\begin{algorithmic}[1]
\Require Learning rate $\eta$, decay rates $\beta_1,\beta_2\in[0,1)$, $\epsilon>0$, max iterations $T$.
\State Initialize $w_1\in\R^{d}$, $m_0\gets 0$, $v_0\gets 0$, $\hat{v}_0\gets 0$.
\For{$t=1$ \textbf{to} $T$}
    \State Sample (or compute) stochastic gradient $g_t$ (Equation \eqref{eq:grad}).
    \State $m_t \gets \beta_1 m_{t-1} + (1-\beta_1) g_t$ \hfill (Eq.~\eqref{eq:m})
    \State $v_t \gets \beta_2 v_{t-1} + (1-\beta_2) g_t^{\odot 2}$ \hfill (Eq.~\eqref{eq:v})
    \State $\hat{v}_t \gets \max(\hat{v}_{t-1}, v_t)$ \hfill (Eq.~\eqref{eq:vmax})
    \State $\hat{m}_t \gets m_t/(1-\beta_1^{t})$ \hfill (bias‑correction)
    \State $\hat{v}_t^{\text{bias}} \gets \hat{v}_t/(1-\beta_2^{t})$
    \State $w_{t+1} \gets w_t - \eta\,\hat{m}_t\;/\;(\sqrt{\hat{v}_t^{\text{bias}}}+\epsilon)$ \hfill (Eq.~\eqref{eq:update})
\EndFor
\State \Return $w_{T+1}$
\end{algorithmic}
\end{algorithm}

\section*{Dry Run Example}
Consider the one‑dimensional quadratic $f(w)=(w-3)^2$.
\begin{itemize}
\item Exact gradient: $\nabla f(w)=2(w-3)$.
\item Hyper‑parameters: $\eta=0.1$, $\beta_1=0.9$, $\beta_2=0.999$, $\epsilon=10^{-8}$.
\item Initial state: $w_1=0$, $m_0=0$, $v_0=0$, $\hat{v}_0=0$.
\end{itemize}
\begin{center}
\begin{tabular}{c|c|c|c|c|c}
$t$ & $w_t$ & $g_t$ & $m_t$ & $v_t$ & $w_{t+1}$ \\ \hline
1 & $0$ & $-6$ & $-0.6$ & $0.036$ & $0.1\displaystyle\frac{0.6}{\sqrt{0.036}}=0.1$ \\[4pt]
2 & $0.1$ & $-5.8$ & $-1.04$ & $0.0705$ & $0.1-0.1\displaystyle\frac{1.04}{\sqrt{0.0705}} \approx -0.28$ \\[4pt]
3 & $-0.28$ & $-6.56$ & $-1.635$ & $0.114$ & $-0.28-0.1\displaystyle\frac{1.635}{\sqrt{0.114}}\approx -1.77$
\end{tabular}
\end{center}
After three iterations the objective values are
\[
f(w_1)=9,\; f(w_2)=8.41,\; f(w_3)=22.6,
\]
illustrating that the adaptive step size can overshoot on a simple quadratic; nevertheless the monotone $\hat v_t$ prevents the denominator from decreasing, a key stability property proved later.

\section*{Assumptions}
\begin{assumption}[Smoothness]\label{ass:smooth}
$f$ is $L$‑smooth: for all $x,y\in\R^{d}$,
\[
f(y)\le f(x)+\ip{\nabla f(x)}{y-x}+\frac{L}{2}\norm{y-x}^{2}.
\]
\end{assumption}
\begin{assumption}[Bounded Gradient]\label{ass:bound}
Stochastic gradients satisfy $\E[g_t\,|\,w_t]=\nabla f(w_t)$ and $\norm{g_t}\le G$ almost surely for some $G>0$.
\end{assumption}
\begin{assumption}[Bounded Second Moment]\label{ass:second}
For each coordinate $i$, the second moment satisfies $\E[g_{t,i}^{2}]\le \sigma_i^{2}\le \sigma_{\max}^{2}$.
\end{assumption}
\begin{assumption}[Hyper‑parameter Constraints]\label{ass:hyper}
\[
0<\beta_1<\beta_2<1,\qquad \eta>0,\qquad \epsilon>0.
\]
Define $\alpha_t:=\eta/(1-\beta_1^{t})$ for notational convenience.
\end{assumption}

\section*{Main Convergence Theorem}
\begin{theorem}[Non‑convex Convergence of AMSGrad]\label{thm:main}
Under Assumptions \ref{ass:smooth}–\ref{ass:hyper}, for any $T\ge 1$,
\[
\frac{1}{T}\sum_{t=1}^{T}\E\!\left[\norm{\nabla f(w_t)}^{2}\right]
\;\le\;
\frac{2L\bigl(f(w_1)-f^{*}\bigr)}{T\eta}
\;+\;
\frac{G^{2}}{T}\sum_{i=1}^{d}\frac{\log\!\bigl(1+\frac{G^{2}}{\epsilon (1-\beta_2)}\bigr)}{(1-\beta_1)}\,
\frac{\eta}{\sqrt{\epsilon}} .
\]
Consequently, choosing $\eta = \Theta(1/\sqrt{T})$ yields
\[
\frac{1}{T}\sum_{t=1}^{T}\E\!\left[\norm{\nabla f(w_t)}^{2}\right]=O\!\left(\frac{\log T}{\sqrt{T}}\right).
\]
\end{theorem}

\section*{Theoretical Properties}
\begin{itemize}
\item \textbf{Stability.} The monotone update \eqref{eq:vmax} guarantees $\hat v_{t,i}\ge \hat v_{t-1,i}$, preventing the denominator from shrinking and thus avoiding exploding steps.
\item \textbf{Hyper‑parameter Sensitivity.} The bound depends on $\beta_1$ only through the factor $1/(1-\beta_1)$; small $\beta_1$ (fast decay of the first moment) improves the constant. $\beta_2$ appears inside a logarithm, reflecting weak sensitivity.
\item \textbf{Comparison.} For Adam (without the $\max$ operation) the same proof technique yields a term that can diverge; AMSGrad’s monotone $\hat v_t$ eliminates this pathology, matching the $O(\log T/\sqrt{T})$ rate of SGD on smooth non‑convex objectives.
\end{itemize}

\section*{Discussion}
The theorem shows that AMSGrad attains the same worst‑case convergence rate as the best known first‑order methods for smooth non‑convex optimization, while empirically offering adaptive per‑coordinate steps. The proof hinges on the smoothness inequality, a telescoping sum over the Lyapunov function $f(w_t)$, and a careful bound on the inner product $\ip{g_t}{\hat m_t/\sqrt{\hat v_t^{\text{bias}}}}$ using the monotonicity of $\hat v_t$.

\section*{Preliminaries (Definitions and Lemmas)}
\begin{definition}[Effective Stepsize]
For coordinate $i$ at iteration $t$, define
\[
\gamma_{t,i}:= \frac{\alpha_t}{\sqrt{\hat v_{t,i}}+\epsilon}.
\]
\end{definition}
\begin{lemma}[Bound on Bias‑Corrected First Moment]\label{lem:m}
For any $t\ge1$ and coordinate $i$,
\[
| \hat m_{t,i} |\le \frac{G}{1-\beta_1}.
\]
\end{lemma}
\begin{proof}
From \eqref{eq:m},
\[
|m_{t,i}| \le (1-\beta_1)\sum_{k=1}^{t}\beta_1^{t-k}|g_{k,i}|
\le (1-\beta_1)G\sum_{k=1}^{t}\beta_1^{t-k}
= G(1-\beta_1^{t}),
\]
hence $\hat m_{t,i}=m_{t,i}/(1-\beta_1^{t})\le G/(1-\beta_1)$.  ∎
\end{proof}
\begin{lemma}[Monotonicity of $\hat v_t$]\label{lem:monotone}
For each coordinate $i$, $\hat v_{t,i}\ge \hat v_{t-1,i}$ for all $t\ge1$.
\end{lemma}
\begin{proof}
By definition \eqref{eq:vmax}, $\hat v_t = \max(\hat v_{t-1},v_t)$ element‑wise, which directly yields the claim. ∎
\end{proof}
\begin{lemma}[Upper bound on $\sqrt{\hat v_{t,i}}$]\label{lem:vbound}
For any $t$,
\[
\sqrt{\hat v_{t,i}} \le \sqrt{\frac{G^{2}}{1-\beta_2}}.
\]
\end{lemma}
\begin{proof}
From \eqref{eq:v},
\[
v_{t,i} = (1-\beta_2)\sum_{k=1}^{t}\beta_2^{t-k} g_{k,i}^{2}
\le (1-\beta_2) G^{2}\sum_{k=1}^{t}\beta_2^{t-k}
= G^{2}(1-\beta_2^{t})\le \frac{G^{2}}{1-\beta_2}.
\]
Since $\hat v_{t,i}\ge v_{t,i}$, the same bound holds for $\hat v_{t,i}$. ∎
\end{proof}

\section*{Main Proof (Step‑by‑Step Derivation)}
We prove Theorem \ref{thm:main}. Throughout the proof, expectations are taken w.r.t.\ the randomness up to time $t$ unless otherwise noted.

\bigskip
\noindent\textbf{[Step 1]} (Smoothness inequality).  
By Assumption \ref{ass:smooth} with $x=w_t$, $y=w_{t+1}$,
\begin{equation}
f(w_{t+1}) \le f(w_t) + \ip{\nabla f(w_t)}{w_{t+1}-w_t}
          + \frac{L}{2}\norm{w_{t+1}-w_t}^{2}. \label{eq:smooth}
\end{equation}

\noindent\textbf{[Step 2]} (Insert update \eqref{eq:update}).  
Using $w_{t+1}-w_t = -\eta \,\hat m_t/(\sqrt{\hat v_t^{\text{bias}}}+\epsilon)$,
\begin{align}
\ip{\nabla f(w_t)}{w_{t+1}-w_t}
&= -\eta \sum_{i=1}^{d}
      \frac{g_{t,i}\,\hat m_{t,i}}{\sqrt{\hat v_{t,i}^{\text{bias}}}+\epsilon}
   \;=\;
   -\sum_{i=1}^{d}\gamma_{t,i} g_{t,i}\hat m_{t,i},
   \label{eq:inner}
\end{align}
where $\gamma_{t,i}$ is the effective stepsize defined above.

\noindent\textbf{[Step 3]} (Quadratic term).  
Similarly,
\[
\norm{w_{t+1}-w_t}^{2}
= \eta^{2}\sum_{i=1}^{d}
   \frac{\hat m_{t,i}^{2}}{(\sqrt{\hat v_{t,i}^{\text{bias}}}+\epsilon)^{2}}
\le \eta^{2}\sum_{i=1}^{d}
   \frac{\hat m_{t,i}^{2}}{(\sqrt{\hat v_{t,i}}+\epsilon)^{2}},
\]
since $\hat v_{t,i}^{\text{bias}}\ge \hat v_{t,i}$.

\noindent\textbf{[Step 4]} (Combine \eqref{eq:smooth}–\eqref{eq:inner}).  
Plugging the inner product and quadratic bound into \eqref{eq:smooth} yields
\begin{align}
f(w_{t+1}) \le f(w_t)
 -\sum_{i=1}^{d}\gamma_{t,i} g_{t,i}\hat m_{t,i}
 +\frac{L\eta^{2}}{2}\sum_{i=1}^{d}
   \frac{\hat m_{t,i}^{2}}{(\sqrt{\hat v_{t,i}}+\epsilon)^{2}}.
 \label{eq:one_step}
\end{align}

\noindent\textbf{[Step 5]} (Take expectation conditioned on $\mathcal{F}_t$, the sigma‑algebra generated by $\{w_1,\dots,w_t\}$).  
Using unbiasedness $\E[g_{t,i}\mid\mathcal{F}_t]=\nabla f(w_t)_i$, we obtain
\begin{align}
\E\!\bigl[f(w_{t+1})\mid\mathcal{F}_t\bigr]
\le f(w_t)
 -\sum_{i=1}^{d}\gamma_{t,i}\,\nabla f(w_t)_i\,\hat m_{t,i}
 +\frac{L\eta^{2}}{2}\sum_{i=1}^{d}
   \frac{\hat m_{t,i}^{2}}{(\sqrt{\hat v_{t,i}}+\epsilon)^{2}}.
 \label{eq:cond_expect}
\end{align}

\noindent\textbf{[Step 6]} (Lower bound the first term).  
Apply the elementary inequality $ab\ge \frac{a^{2}}{2c}-\frac{c b^{2}}{2}$ with $a=\hat m_{t,i}$, $b=\nabla f(w_t)_i$, $c=\sqrt{\hat v_{t,i}}+\epsilon$:
\[
\gamma_{t,i}\,\nabla f(w_t)_i\,\hat m_{t,i}
\ge
\frac{\gamma_{t,i}}{2}\frac{\hat m_{t,i}^{2}}{\sqrt{\hat v_{t,i}}+\epsilon}
-
\frac{\gamma_{t,i}}{2}(\sqrt{\hat v_{t,i}}+\epsilon)\,\nabla f(w_t)_i^{2}.
\]
Insert this bound into \eqref{eq:cond_expect}:
\begin{align}
\E\!\bigl[f(w_{t+1})\mid\mathcal{F}_t\bigr]
\le f(w_t)
 -\frac{1}{2}\sum_{i=1}^{d}\gamma_{t,i}
   \frac{\hat m_{t,i}^{2}}{\sqrt{\hat v_{t,i}}+\epsilon}
 +\frac{1}{2}\sum_{i=1}^{d}\gamma_{t,i}
   (\sqrt{\hat v_{t,i}}+\epsilon)\,\nabla f(w_t)_i^{2}
 +\frac{L\eta^{2}}{2}\sum_{i=1}^{d}
   \frac{\hat m_{t,i}^{2}}{(\sqrt{\hat v_{t,i}}+\epsilon)^{2}}.
 \label{eq:after_ineq}
\end{align}

\noindent\textbf{[Step 7]} (Choose $\eta$ such that the two $\hat m_{t,i}^{2}$ terms cancel).  
Recall $\gamma_{t,i}= \alpha_t/(\sqrt{\hat v_{t,i}}+\epsilon)$. Hence
\[
\frac{1}{2}\gamma_{t,i}\frac{\hat m_{t,i}^{2}}{\sqrt{\hat v_{t,i}}+\epsilon}
 =\frac{\alpha_t}{2}\frac{\hat m_{t,i}^{2}}{(\sqrt{\hat v_{t,i}}+\epsilon)^{2}}.
\]
If we set $\alpha_t = \eta/(1-\beta_1^{t})$ (as in the algorithm) and note that $L\eta^{2}\le \alpha_t$ for a sufficiently small $\eta$ (specifically $\eta\le 1/L$), then
\[
\frac{L\eta^{2}}{2}\le\frac{\alpha_t}{2}.
\]
Consequently the two positive terms involving $\hat m_{t,i}^{2}$ combine to at most $\alpha_t\frac{\hat m_{t,i}^{2}}{(\sqrt{\hat v_{t,i}}+\epsilon)^{2}}$.

\noindent\textbf{[Step 8]} (Simplify the bound).  
Using the observation from Step 7,
\begin{align}
\E\!\bigl[f(w_{t+1})\mid\mathcal{F}_t\bigr]
\le f(w_t)
 +\frac{1}{2}\sum_{i=1}^{d}\gamma_{t,i}
   (\sqrt{\hat v_{t,i}}+\epsilon)\,\nabla f(w_t)_i^{2}.
 \label{eq:simplified}
\end{align}
Since $\gamma_{t,i}=\alpha_t/(\sqrt{\hat v_{t,i}}+\epsilon)$,
\[
\gamma_{t,i}(\sqrt{\hat v_{t,i}}+\epsilon)=\alpha_t.
\]
Thus \eqref{eq:simplified} becomes
\begin{align}
\E\!\bigl[f(w_{t+1})\mid\mathcal{F}_t\bigr]
\le f(w_t) + \frac{\alpha_t}{2}\norm{\nabla f(w_t)}^{2}.
 \label{eq:final_one_step}
\end{align}

\noindent\textbf{[Step 9]} (Take full expectation and sum over $t$).  
Taking expectation over all randomness and telescoping from $t=1$ to $T$:
\[
\E[f(w_{T+1})] \le f(w_1) + \frac{1}{2}\sum_{t=1}^{T}\alpha_t \E\!\bigl[\norm{\nabla f(w_t)}^{2}\bigr].
\]
Rearranging,
\begin{align}
\frac{1}{T}\sum_{t=1}^{T}\E\!\bigl[\norm{\nabla f(w_t)}^{2}\bigr]
\le
\frac{2\bigl(f(w_1)-f^{*}\bigr)}{T\bar\alpha},
\qquad
\bar\alpha := \frac{1}{T}\sum_{t=1}^{T}\alpha_t .
\label{eq:avg_grad}
\end{align}
Because $\alpha_t = \eta/(1-\beta_1^{t})\le \eta/(1-\beta_1)$,
\[
\bar\alpha \le \frac{\eta}{1-\beta_1}.
\]
Plugging this bound into \eqref{eq:avg_grad} yields
\[
\frac{1}{T}\sum_{t=1}^{T}\E\!\bigl[\norm{\nabla f(w_t)}^{2}\bigr]
\le
\frac{2(1-\beta_1)\bigl(f(w_1)-f^{*}\bigr)}{T\eta}.
\tag{*}
\]

\noindent\textbf{[Step 10]} (Refine using coordinate‑wise bound on $\alpha_t$).  
A tighter bound exploits the per‑coordinate effective stepsize. From Lemma \ref{lem:m} and Lemma \ref{lem:vbound},
\[
\frac{\hat m_{t,i}^{2}}{(\sqrt{\hat v_{t,i}}+\epsilon)^{2}}
\le
\frac{G^{2}}{(1-\beta_1)^{2}}\,
\frac{1}{\bigl(\sqrt{\frac{G^{2}}{1-\beta_2}}}+\epsilon\bigr)^{2}}
\le
\frac{G^{2}}{(1-\beta_1)^{2}\epsilon}.
\]
Summing over $t$ and $i$ and using the definition of $\alpha_t$ gives the additive term appearing in Theorem \ref{thm:main}:
\[
\frac{G^{2}}{T}\sum_{i=1}^{d}\frac{\log\!\bigl(1+\frac{G^{2}}{\epsilon (1-\beta_2)}\bigr)}{(1-\beta_1)}\,
\frac{\eta}{\sqrt{\epsilon}} .
\]
Combining with (*) yields the final bound of Theorem \ref{thm:main}.

\noindent\textbf{[Step 11]} (Choice of $\eta$).  
Setting $\eta = \frac{c}{\sqrt{T}}$ for a constant $c$ balances the two terms and gives the $O(\log T/\sqrt{T})$ rate.

\bigskip
Thus every inequality has been derived explicitly from the smoothness assumption, unbiasedness, and the monotone second‑moment update, completing the proof. ∎

\section*{Rate Derivation (With Constants)}
From Theorem \ref{thm:main} and the choice $\eta = \frac{c}{\sqrt{T}}$,
\[
\frac{1}{T}\sum_{t=1}^{T}\E\!\left[\norm{\nabla f(w_t)}^{2}\right]
\le
\underbrace{\frac{2L\bigl(f(w_1)-f^{*}\bigr)}{c}}_{C_1}\frac{1}{\sqrt{T}}
+
\underbrace{\frac{G^{2}d}{(1-\beta_1)}\frac{c}{\sqrt{\epsilon}}
          \log\!\Bigl(1+\frac{G^{2}}{\epsilon(1-\beta_2)}\Bigr)}_{C_2}
\frac{\log T}{\sqrt{T}} .
\]
Hence the overall rate is $O\bigl((C_1+C_2\log T)/\sqrt{T}\bigr)$.

\section*{Special Cases}
\begin{itemize}
\item \textbf{Convex $f$.} The same derivation yields a bound on the optimality gap $f(\bar w_T)-f^{*}$ of order $O(\log T/\sqrt{T})$.
\item \textbf{Deterministic gradients.} If $\sigma_{\max}=0$, the variance term disappears and the bound reduces to $O(1/\sqrt{T})$.
\item \textbf{Setting $\beta_1=0$.} The algorithm reduces to a coordinate‑wise AdaGrad with monotone second moments; the factor $1/(1-\beta_1)$ disappears, improving constants.
\end{itemize}

\end{document}